% New slide.

\begin{frame}[t, fragile]{Using cabal to manage and build Haskell projects.}

Haskell has its own project management and build tool. This tool is called
cabal and it is a command line utility which is designed to operate on
files that have a .cabal filename extension.

However, before cabal can be used to manage and build a particular Haskell
project, it must first initialise that project by running the following
command from within the project's top-level directory; \\

\begin{minted}[frame=single]{console}
$ cabal init
\end{minted}

\end{frame}


% New slide.

\begin{frame}[t, fragile]{Using .cabal files to build a binary executable.}

An example of a simple .cabal file is shown on the next slide.

This file contains just the one target. As you can probably guess
from the use of the keyword ``executable'', this target can be
used to build a binary executable. The name of the target is
basicExample and it is defined on lines 10-18 of the cabal file.
Incidentally, the name of the target is assigned to the binary
executable that results from successfully invoking this target.

To build the binary executable in this cabal file, invoke the
following command from the command line;

\begin{minted}[frame=single]{console}
$ cabal exe:basicExample
\end{minted}

\end{frame}